\documentclass[12pt]{article}
%---DOCUMENT MARGINS---
\usepackage{geometry} % Required for adjusting page dimensions and margins
\geometry{
	paper=a4paper, % Paper size, change to letterpaper for US letter size
	top=4cm, % Top margin
	bottom=2cm, % Bottom margin
	left=2cm, % Left margin
	right=2cm, % Right margin
	headheight=3cm, % Header height
	%footskip=1.5cm, % Space from the bottom margin to the baseline of the footer
	headsep=0.5cm, % Space from the top margin to the baseline of the header
	%showframe, % Uncomment to show how the type block is set on the page
}
\usepackage{cmap}

\usepackage[T1,T2A]{fontenc}
\usepackage{fontspec}
\setmainfont[
	BoldFont={* Bold},
	ItalicFont={* Italic},
	BoldItalicFont={* BoldItalic}
]{DejaVu Serif Condensed}
\setsansfont{DejaVu Sans}
\setmonofont{Dejavu Sans Mono}
\usepackage[math-style=TeX]{unicode-math}
\setmathfont{Latin Modern Math}

\usepackage[nobottomtitles*]{titlesec}
\titleformat
{\section} % command
{\fontsize{14}{14}\sffamily\bfseries} % format
{} % label
{0pt} % sep
{} % before-code
\titlespacing{\section}{0pt}{0.75em}{0.25em}
\newcommand{\tl}{}
\newcommand{\ml}{}
\newcommand{\problem}[1]{%
	\section[#1]{#1 \fontsize{14}{14}\selectfont{\hfill{\normalfont{
				\emoji{hourglass-not-done}\tl \space\space \emoji{floppy-disk} \ml}}}}
}
\AddToHook{cmd/section/before}{\clearpage}

\usepackage[dvipsnames]{xcolor}
\titleformat
{\subsection} % command
{\fontsize{14}{14}\sffamily\bfseries} % format
{\includesvg[width=0.035\textwidth, inkscapelatex=false]{lion}\space} % label
{0pt} % sep
{} % before-code
\titlespacing{\subsection}{0pt}{0.75em}{0.25em}

\setlength{\parskip}{0.5em}
\setlength{\parindent}{0pt}
\renewcommand{\baselinestretch}{1.2}
\sloppy

\usepackage{listings}
\definecolor{lstbg}{RGB}{250,250,252}
\definecolor{lstframe}{RGB}{220,220,225}
\lstdefinestyle{mystyle}{
	backgroundcolor=\color{lstbg},
	frame=single,
	rulecolor=\color{lstframe},
	basicstyle=\ttfamily\small,
	keywordstyle=\bfseries,
	breaklines=true,
	aboveskip=0.5\baselineskip,
	belowskip=0\baselineskip  
}
\lstset{style=mystyle, language=C++}

\usepackage{fancyhdr}
\pagestyle{fancy}
\setlength{\headwidth}{\textwidth}
\newcommand{\round}{}
\newcommand{\problemName}{}
\newcommand{\lang}{English}
\fancyhead[L]{
	\begin{minipage}{0.4\textwidth}
		\bf{
			EJOI 2025 \round\\
			\problemName\\
			\emoji{globe-with-meridians} \lang
		}
	\end{minipage}
	\vspace{0.5cm}
}
\fancyhead[R]{%
	\begin{minipage}{0.6\textwidth}
		\includesvg[width=\textwidth, inkscapelatex=false]{logo}
	\end{minipage}
	\vspace{0.5cm}
}
\usepackage{lastpage}
\fancyfoot[C]{\problemName\space(\lang) \thepage\ / \pageref*{LastPage}}

\raggedbottom

\usepackage{amsmath}
\usepackage{stmaryrd}

\usepackage{graphicx}
\graphicspath{{./}}
\usepackage[export]{adjustbox}
\usepackage{wrapfig}
\makeatletter
\patchcmd\WF@putfigmaybe{\lower\intextsep}{}{}{\fail}
\AddToHook{env/wrapfigure/begin}{\setlength{\intextsep}{0pt}}
\makeatother
\usepackage[inkscapearea=page, inkscapepath=./svg-inkscape]{svg}
\svgpath{{./}}

\usepackage{makecell}
\usepackage{tabularray}
\AtBeginEnvironment{table}{\vspace{-0.2cm}}
\AtEndEnvironment{table}{\vspace{-0.2cm}}
\usepackage{float}
\usepackage{placeins}
\usepackage{caption}
\captionsetup[table]{
	skip=0.25em,
	singlelinecheck=false, justification=justified
}

\usepackage{enumitem}
\setlist{itemsep=-0.4em, leftmargin=22pt, topsep=-\parskip}
\AfterEndEnvironment{itemize}{\vspace{\parskip}}
\newcommand{\tabitem}{\indent~~\llap{\textbullet}~~}

\usepackage{hyperref}
\hypersetup{
	colorlinks=true,
	citecolor=blue,
	linkcolor=blue,
	urlcolor=cyan,
}

\usepackage{emoji}

\usepackage[english]{babel}

\begin{document}

\renewcommand{\round}{Day 1}
\renewcommand{\problemName}{Task Prison}
\renewcommand{\lang}{Hungarian [HUN]}

\renewcommand{\tl}{$2$ sec.}
\renewcommand{\ml}{$1024$ MB}
\problem{Börtön}

% Alice and Bob have been unjustly sentenced to a maximum-security prison. Now they must plan their escape. To do this, they need to be able to communicate as efficiently as possible (in particular, Alice needs to send daily information to Bob). However, they cannot meet up and can only exchange information through notes written on napkins. Each day Alice wants to send a new piece of information to Bob -- a number between $0$ and $N-1$. At every lunch, Alice gets three napkins and writes a number between $0$ and $M-1$ on each napkin (there may be repetitions) and leaves them on her seat. Then, their enemy, Charly, destroys one of the napkins and mixes the other two up. Finally, Bob finds the two remaining napkins and reads the numbers on them. He must accurately decode the original number that Alice wanted to send him. There is limited space on the napkins, so $M$ is fixed. However, Alice and Bob's goal is to maximize the information throughput, so they are free to choose $N$ as large as they can. Help Alice and Bob by implementing a strategy for each of them, trying to maximize the value of $N$.
Alice-t és Bobot igazságtalanul bezárták egy szigorúan őrzött börtönbe. Most a szökésüket tervezik. Ehhez a lehető leghatékonyabban kell kommunikálniuk (pontosabban, Alice minden nap információt szeretne küldeni Bobnak). Sajnos nem tudnak találkozni, így csak szalvétákra írt számokkal tudnak kommunikálni. Alice minden nap egy új információt szeretne küldeni Bobnak -- egy számot $0$ és $N-1$ között. Minden délben Alice fog három szalvétát és mindháromra ír egy-egy számot $0$ és $M-1$ között (a számok lehetnek egyformák), majd otthagyja őket egy széken. Ezután az ellenségük, Charly megsemmisíti az egyik szalvétát és összekeveri a maradék kettőt. Végül Bob elolvassa a maradék két szalvétára írt számokat. Bobnak pontosan vissza kell állítania az eredeti számot, amit Alice küld neki. A szalvétákon kevés a hely, így $M$ rögzített. Alice és Bob célja, hogy maximalizálják az átadható információ mennyiségét, így a lehető legnagyobb $N$-t választhatják. Segíts Alice-nak és Bobnak azzal, hogy megvalósítasz egy stratégiát mindkettejük számára, mely maximalizálja $N$ értékét.

%\subsection{Implementation details}
\subsection{Megvalósítás}
% Since this is a communication problem, your program will be run in two separate executions (one for Alice and one for Bob) that cannot share data or communicate in any way other than the one described here. You need to implement three functions:
Mivel ez egy kommunkációs probléma, a kódodat két részletben futtatjuk (egyszer mint Alice és egyszer mint Bob). A futtatások nem oszthatnak meg egymással adatot és nem kommunikálhatnak semmilyen más módon, az eddig ismertetett módszeren kívül. Három függvényt kell megvalósítanod:

\begin{lstlisting}
 int setup(int M);
\end{lstlisting}

% This will be called once at the start of Alice's execution of your program and once at the start of Bob's execution. It is given $M$ and must return the desired $N$. Both calls to \texttt{setup} must return the same $N$.
Ezt a függvényt egyszer hívja meg a program Alice futtatásának kezdetén, egyszer pedig Bob futtatásának kezdetén. A függvény bemenete $M$, és a kívánt $N$-t kell visszaadnia. A \texttt{setup} függvény mindkét hívásának ugyanazt az $N$-t kell visszaadnia.

\begin{lstlisting}
 std::vector<int> encode(int A);
\end{lstlisting}

% This implements Alice's strategy. It will be called with the number to encode $A$ ($0 \leq A < N$) and must return three numbers $W_1, W_2, W_3$ ($0 \leq W_i < M$) that encode $A$. This function will be called a total of $T$ times -- once per day (values of $A$ may repeat between days).
Ez valósítja meg Alice stratégiáját. Bemenete az elkódolandó szám $A$ ($0 \leq A < N$) és három számot kell visszaadnia $W_1, W_2, W_3$ ($0 \leq W_i < M$) melyek az $A$ számot kódolják. Ez a függvény $T$ alkalommal lesz meghívva -- minden naphoz egyszer (ugyanaz $A$ érték több napon is előfordulhat).

\begin{lstlisting}
 int decode(int X, int Y);
\end{lstlisting}

% This implements Bob's strategy. It will be called with two of the three numbers returned by \texttt{encode} in some order. It must return the same value $A$ that encode received. This function will also be called $T$ times -- corresponding to the $T$ calls to \texttt{encode}; they will be in the same order. All calls to \texttt{encode} will happen before all calls to \texttt{decode}.
Ez valósítja meg Bob stratégiáját. A bemenete az \texttt{encode} függvény által előállított három szám közül kettő, valamilyen sorrendben. Azt az $A$ számot kell visszaadnia, amelyiket az \texttt{encode} megkapta. Ez a függvény is $T$-szer lesz meghívva -- az \texttt{encode} mind a $T$ hívásához egyszer, azonos sorrendben. Minden \texttt{encode} hívás előbb történik meg, mint bármelyik \texttt{decode} hívás.

%\subsection{Constraints}
\subsection{Korlátok}

\begin{itemize}
\item $M \leq 4300$
\item $T = 5000$
\end{itemize}

%\subsection{Scoring}
\subsection{Pontozás}

%For a particular subtask, the fraction S of the points you get depends on the smallest $N$ returned by \texttt{setup} on any test in that subtask. It also depends on $N^*$, which is the target value of $N$ that you need to get the full points for the subtask:
Egy adott részfeladat esetében a kapott pontok S hányada attól függ, hogy mi a legkisebb $N$, amit a \texttt{setup} függvény az adott részfeladat bármelyik tesztjénél visszaad. Függ továbbá az $N^*$ értéktől is, amely az $N$ célértéke, amire a részfeladat teljes pontszámának megszerzéséhez szükség van:

\begin{itemize}
%\item If your solution fails on any test, then $S = 0$.
\item Ha a megoldásod bármelyik teszten rossz eredményt ad, akkor $S = 0$.
\item Ha $N \geq N^*$, akkor $S = 1.0$.
\item Ha $N < N^*$, akkor $S = \max\left(0.35 \max\left(\frac{\log(N) - 0.985 \log(M)}{\log(N^*) - 0.985 \log(M)}, 0.0\right)^{0.3} + 0.65 \left(\frac{N}{N^*}\right)^{2.4}, 0.01\right)$.
\end{itemize}

%\subsection{Subtasks}
\subsection{Részfeladatok}

\begin{table}[H]
\begin{tblr}{|Q[c,m]|Q[c,m]|Q[c,m]|Q[c,m]|}
\hline
%\textbf{Subtask} & \textbf{Points} & $M$ & $N^*$ \\
\textbf{Részfeladat} & \textbf{Pontszám} & $M$ & $N^*$ \\
\hline
$1$ & $10$ & $700$ & $82017$ \\
\hline
$2$ & $10$ & $1100$ & $202217$ \\
\hline
$3$ & $10$ & $1500$ & $375751$ \\
\hline
$4$ & $10$ & $1900$ & $602617$ \\
\hline
$5$ & $10$ & $2300$ & $882817$ \\
\hline
$6$ & $10$ & $2700$ & $1216351$ \\
\hline
$7$ & $10$ & $3100$ & $1603217$ \\
\hline
$8$ & $10$ & $3500$ & $2043417$ \\
\hline
$9$ & $10$ & $3900$ & $2536951$ \\
\hline
$10$ & $10$ & $4300$ & $3083817$ \\
\hline
\end{tblr}
\end{table}
\FloatBarrier

%\subsection{Example}
\subsection{Példa}

% Consider the following example with $T = 5$. We consider an encoding scheme where Alice sends three equal numbers to encode $0$ or three distinct numbers to encode $1$. Notice that Bob can decode the original number from any two of the three numbers Alice sent.
Tekintsük a következő példát $T = 5$ értékkel. Tekintsünk egy olyan kódolási sémát, amelyben Alice három egyenlő számot küld a $0$ kódolásához, vagy három különböző számot az $1$ kódolásához. Vegyük észre, hogy Bob az eredeti számot az Alice által küldött három szám közül bármelyikből kettőből dekódolni tudja.

\begin{table}[H]
\begin{tblr}{|Q[c,m]|Q[c,m]|Q[c,m]|}
\hline
%\textbf{\makecell{Execution}} & \textbf{\makecell{Function call}} & \textbf{Return value} \\
\textbf{\makecell{Futtatás}} & \textbf{\makecell{Függvényhívás}} & \textbf{Visszatérési érték} \\
\hline
Alice & \texttt{setup(10)} & \texttt{2} \\
\hline
Bob & \texttt{setup(10)} & \texttt{2} \\
\hline
Alice & \texttt{encode(0)} & \texttt{\{5, 5, 5\}} \\
\hline
Alice & \texttt{encode(1)} & \texttt{\{8, 3, 7\}} \\
\hline
Alice & \texttt{encode(1)} & \texttt{\{0, 3, 1\}} \\
\hline
Alice & \texttt{encode(0)} & \texttt{\{7, 7, 7\}} \\
\hline
Alice & \texttt{encode(1)} & \texttt{\{6, 2, 0\}} \\
\hline
Bob & \texttt{decode(5, 5)} & \texttt{0} \\
\hline
Bob & \texttt{decode(8, 7)} & \texttt{1} \\
\hline
Bob & \texttt{decode(3, 0)} & \texttt{1} \\
\hline
Bob & \texttt{decode(7, 7)} & \texttt{0} \\
\hline
Bob & \texttt{decode(2, 0)} & \texttt{1} \\
\hline
\end{tblr}
\end{table}

%\subsection{Sample grader}
\subsection{Mintaértékelő}

% For the sample grader, all calls to \texttt{encode} and \texttt{decode} will be in the same execution of your program. Additionally, \texttt{setup} will be called only once (as opposed to twice, once per execution, as in the grading system).
A mintaértékelőben minden \texttt{encode} és \texttt{decode} hívás egy programfuttatásban történik. Ebben a \texttt{setup} függvény csak egyszer kerül meghívásra (ellentétben az értékelő rendszerrel, ahol egyszer-egyszer mindkét program futtatásakor).

% The input is just a single integer -- $M$. Then it will print out the $N$ your \texttt{setup} returned. It will then \texttt{encode} and \texttt{decode} $T$ times with randomly generated numbers from $0$ to $N-1$ and randomly generated choices of which two of the three numbers from \texttt{encode} to give to \texttt{decode} (and in what order). It will print out an error message if your solution failed.
A bemenet egyetlen egész szám -- $M$. A mintaértékelő először kiírja az $N$ értéket, amit a \texttt{setup} visszaadott. Ezután $T$ alkalommal meghívja az \texttt{encode} és \texttt{decode} függvényeket véletlenszerűen generált számokkal $0$ és $N-1$ között, és véletlenszerűen generált választásokkal arra vonatkozóan, hogy az \texttt{encode} által generált három szám közül melyik kettőt adja a \texttt{decode}-nak (és milyen sorrendben). A program hibaüzenetet ír ki, ha a megoldásod hibás.

\end{document}
